\documentclass[11pt]{article}
\usepackage[T1]{fontenc}
\usepackage[margin=1in]{geometry}
\usepackage{fancyvrb}
\usepackage{xcolor}
\usepackage{enumitem}
\usepackage[hidelinks]{hyperref}

\definecolor{HardyBlue}{RGB}{24,64,112}
\setlength{\parindent}{0pt}
\setlength{\parskip}{0.7em}
\setlist[itemize]{leftmargin=1.5em}

\begin{document}

\begin{center}
{\LARGE\color{HardyBlue}\bfseries The Irrationality of √2 + √3\par}
\vspace{1em}
Formal status: Kernel verified\\
Faithfulness: User approved; independently reviewed by claude-opus-5 on claude (agreed)\\
Informal completeness: Not independently assessed\\
Document status: TeX compiled
\end{center}

\hrule
\vspace{1em}

\section*{Theorem}
Theorem. The real number √2 + √3 is irrational; that is, there is no rational number q with (q : ℝ) = √2 + √3.

Lean statement (Frozen Claim 2c1618098577b6a91f8602cd4e53d9ac2a361f979ca93444bd3ce32d35c4a06b):

    theorem irrational\_sqrt\_two\_add\_sqrt\_three : Irrational (Real.sqrt (2 : ℝ) + Real.sqrt (3 : ℝ))

Here `Irrational x` abbreviates `x ∉ Set.range ((↑) : ℚ → ℝ)`, so the statement is exactly the assertion that no rational number, viewed as a real, equals √2 + √3.

\section*{Approved interpretation}
\begin{itemize}
\item The real numbers ℝ (Lean/Mathlib's `Real`), a complete ordered field
\item The rational numbers ℚ, embedded in ℝ by the coercion `((↑) : ℚ → ℝ)`
\item `Real.sqrt : ℝ → ℝ`, the nonnegative square-root function on ℝ (returning 0 on negative inputs)
\item The user's claim is closed: it asserts a property of one fixed real number, so the theorem statement has no explicit binders.
\item One implicit universal quantifier is hidden inside `Irrational`: `Irrational x` unfolds to `x ∉ Set.range ((↑) : ℚ → ℝ)`, i.e. `∀ q : ℚ, (q : ℝ) ≠ x`.
\item None; the claim is unconditional, so the theorem has no hypotheses.
\item Standard Mathlib meanings for `Real`, `Real.sqrt`, and `Irrational` are taken as given.
\item The arguments 2 and 3 are the real numbers 2 and 3; being nonnegative, `Real.sqrt 2` and `Real.sqrt 3` are the genuine positive square roots and not the junk value 0.
\item 'sqrt(2)' and 'sqrt(3)' are the principal (nonnegative) real square roots, written `Real.sqrt (2 : ℝ)` and `Real.sqrt (3 : ℝ)`.
\item 'is irrational' is Mathlib's `Irrational : ℝ → Prop`, defined as `x ∉ Set.range ((↑) : ℚ → ℝ)`. This is the standard reading; it neither weakens nor strengthens the claim.
\item Numerals are given explicit type ascriptions `(2 : ℝ)` and `(3 : ℝ)` so that elaboration fixes their type as ℝ rather than leaving a metavariable for the default-instance mechanism; this is a purely syntactic repair and does not change the mathematical content.
\item Addition is the field addition of ℝ.
\item No stronger algebraic claim (e.g. that √2 + √3 has degree 4 over ℚ, or that √2 and √3 are individually irrational) is asserted, since the user's claim is only about irrationality of the sum.
\end{itemize}

\section*{Human-readable proof}
Overview. The proof squares the number in question. Squaring turns the sum of two surds into a rational shift of a single surd, √6, and the irrationality of √6 is a standard fact available in Mathlib. Formally the argument is a short reduction: rationality of √2 + √3 would force rationality of √6.

Step 1: √6 is irrational.
The integer 6 is not a perfect square, since 2² = 4 < 6 < 9 = 3². Mathlib's criterion `irrational\_sqrt\_natCast\_iff` (equivalently `irrational\_sqrt\_ofNat\_iff`) states that for a natural number n, `Irrational (√n) ↔ ¬IsSquare n`. Applying it to n = 6 gives `Irrational (√6)`. In the formal proof this whole step is discharged by `norm\_num`, which carries a dedicated irrationality extension (`Mathlib.Tactic.NormNum.Irrational`) implementing exactly this criterion; a direct `decide` on `¬IsSquare 6` does not reduce, so the criterion is invoked either through `norm\_num` or with kernel-level decision.

Step 2: elementary facts about the two square roots.
Since 2 ≥ 0 and 3 ≥ 0, the defining property of `Real.sqrt` on nonnegative arguments gives
    (√2)² = 2 and (√3)² = 3,
formally `Real.sq\_sqrt`. Multiplicativity of the square root on nonnegative arguments (`Real.sqrt\_mul`, applied with 0 ≤ 2) gives
    √2 · √3 = √(2 · 3) = √6.

Step 3: the reduction.
Suppose, for contradiction, that √2 + √3 were rational. Unfolding `Irrational`, this means there is a rational q with
    (q : ℝ) = √2 + √3.
Squaring both sides and expanding,
    q² = (√2 + √3)² = (√2)² + 2·√2·√3 + (√3)² = 2 + 2√6 + 3 = 5 + 2√6,
using Step 2 for the three substitutions. Solving for √6,
    √6 = (q² − 5)/2.

The right-hand side is the image in ℝ of the rational number (q² − 5)/2, since ℚ is closed under squaring, subtraction, and division by 2. Hence √6 lies in the range of the coercion ℚ → ℝ, i.e. √6 is rational. This contradicts Step 1.

Therefore no such rational q exists, and √2 + √3 is irrational. ∎

Remarks on the formal rendering. In Lean the contradiction is packaged directly rather than by an explicit `by\_contra`: `rintro ⟨q, hq⟩` introduces the hypothetical rational witness q together with `hq : (q : ℝ) = √2 + √3`, and `refine hirr ⟨(q \textasciicircum{} 2 - 5) / 2, ?\_⟩` supplies the rational witness for √6 that Step 1 forbids. The remaining goal, `((((q \textasciicircum{} 2 - 5) / 2 : ℚ)) : ℝ) = √6`, is closed by moving the coercion inside with `push\_cast` and applying `linarith` against the identity `q² = 5 + 2√6` established from `hq` together with (√2)² = 2, (√3)² = 3, and √2·√3 = √6 (the expansion of the square is handled by `nlinarith`).

Status. The Lean proof body corresponding to this exposition was checked against the Frozen Claim with Hardy's proof-checking tool under the pinned toolchain (Lean 4.33.1, Mathlib revision 0df444a360eaa60ab8c11dca51a86af692955474) and reported no errors and no open goals. That is tool feedback only. This informal exposition has not been independently assessed, and acceptance of the solution rests solely with Hardy's FinalVerifier.

\section*{Known gaps}
\begin{itemize}
\item This prose exposition has not been independently assessed. It is a human-readable rendering of the formal argument, and any discrepancy between the prose and the machine-checked Lean term would not have been detected by the checks performed here.
\item The irrationality of √6 is not proved from first principles in this document. It is imported as a Mathlib fact (`irrational\_sqrt\_natCast\_iff` / `irrational\_sqrt\_ofNat\_iff`, applied via the `norm\_num` irrationality extension), and the exposition only justifies the side condition that 6 is not a perfect square.
\item The proof is relative to the correctness of the pinned Lean and Mathlib versions and of Hardy's checking harness; a successful tool run is feedback, not acceptance. Only Hardy's FinalVerifier can accept the solution.
\item No claim is made here about stronger or related statements — for instance the irrationality of √2 + √3 over larger fields, its algebraic degree (it is in fact algebraic of degree 4, a root of x⁴ − 10x² + 1), or the irrationality of √2 and √3 individually. These are not needed for, and not established by, the argument given.
\end{itemize}

\section*{Exact Lean statement}
\begin{Verbatim}[fontsize=\small]
theorem irrational_sqrt_two_add_sqrt_three : Irrational (Real.sqrt (2 : ℝ) + Real.sqrt (3 : ℝ))
\end{Verbatim}

\section*{Verification appendix}
\textbf{Axioms used:} propext, Classical.choice, Quot.sound

\begin{Verbatim}[fontsize=\small]
Frozen Claim SHA-256: 2c1618098577b6a91f8602cd4e53d9ac2a361f979ca93444bd3ce32d35c4a06b
Lean source SHA-256: e9adc3c703f981860651c4ef8b8ca8540745152a22b483672b3793a857665912
Verification SHA-256: b158ba7d4c8ee2b98bb0aa9fc81044873016df1558df27456339324f5b0a20b1
Prompt set SHA-256: 50c207a708e3696c9846957a60089ccd0fb4a6ddbc6a6920449887455cf3f5e2
Tectonic bundle SHA-256: 6ffe055852f8faf66c0acbe1a7fb27f87b869a90bad1204f3bf4d9683f597c7c
\end{Verbatim}

\begin{samepage}
\textbf{Run and tool identities}
\begin{itemize}
\item Run ID: 0e543e24-002e-4f43-bdc6-5c09a86d95a6
\item Model: claude-opus-5
\item Backend: claude
\item Runtime SDK: 0.2.127
\item Lean: 4.33.1
\item Mathlib: 0df444a360eaa60ab8c11dca51a86af692955474
\item Tectonic: 0.16.9
\item Tectonic bundle: https://data1.fullyjustified.net/tlextras-2022.0r0.tar
\end{itemize}
\end{samepage}

\vfill
\small
Lean kernel verification, faithfulness (approved by the user and read back by an
independent model before proof search), heuristic informal review, and document
compilation are independent statuses. This trusted-development MVP
does not sandbox generated Lean or TeX.

\end{document}
